\begin{figure}[htbp]
\centering
\begin{tikzpicture}[node distance=2.8cm, x=1cm, y=1cm]
    % Glucose box
    \draw[thick, fill=UofSCGarnet!25] (-1,0.8) rectangle (1,1.8);
    \node[font=\large] at (0,1.3) {Glucose $G$};

    % Insulin box
    \draw[thick, fill=UofSCAtlantic!25] (-1,-1.2) rectangle (1,-0.2);
    \node[font=\large] at (0,-0.7) {Insulin $I$};

    % Meal input arrow
    \draw[->, thick, UofSCHorseshoe] (-2.8,1.3) -- (-1,1.3) node[left, pos=0, font=\small] {Meal $D$};

    % Insulin injection arrow
    \draw[->, thick, UofSCHorseshoe] (-2.8,-0.7) -- (-1,-0.7) node[left, pos=0, font=\small] {Injection $u$};

    % Glucose utilization arrow
    \draw[->, thick, UofSCCongaree] (1,1.3) -- (2.8,1.3) node[right, pos=1, font=\small] {Utilization};

    % Insulin clearance arrow
    \draw[->, thick, UofSCCongaree] (1,-0.7) -- (2.8,-0.7) node[right, pos=1, font=\small] {Clearance};

    % Insulin enhances glucose uptake (upward arrow from insulin to glucose)
    \draw[->, thick, UofSCAtlantic] (0.3,-0.2) -- (0.3,0.8) node[midway, right, font=\small] {Enhances};
    \node[right, font=\tiny] at (0.6,0.3) {uptake};

    % Glucose stimulates insulin secretion (downward arrow from glucose to insulin)
    \draw[->, thick, UofSCGarnet] (-0.3,0.8) -- (-0.3,-0.2) node[midway, left, font=\small] {Stimulates};
    \node[left, font=\tiny] at (-0.8,0.3) {secretion};
\end{tikzpicture}
\caption{Glucose-insulin feedback loop block diagram showing meal and infusion inputs, utilization/clearance outputs, and the bidirectional regulatory feedback.}
\label{fig:glucose_insulin}
\end{figure}
